\section{Design an Elevator System}
\label{sec:case-elevator}

\problemstatement{Design the control system for a building with multiple
elevators. Handle external hall calls (up/down on a floor) and internal
car requests (go to floor N), decide which elevator serves a request, and
move each car efficiently.}

\begin{patternbox}[Clarify Before Designing]
Ask: How many elevators and floors? What scheduling goal -- minimise wait,
minimise travel, or fairness? Are there express/service elevators? For
this design assume: $N$ elevators, a central dispatcher, a
nearest-suitable-car scheduling policy that is swappable, and each car
running a SCAN (elevator algorithm) sweep.
\end{patternbox}

\subsection*{Identify the entities}

\texttt{Elevator} (a car with current floor, direction, and its set of
target floors), \texttt{Request} (hall or car call), \texttt{Direction}
(up/down/idle), a \texttt{Dispatcher} that assigns requests to cars via a
\texttt{SchedulingStrategy}, and a \texttt{Building} owning the elevators.
Each elevator behaves as a small state machine: idle, moving up, moving
down.

\begin{ideabox}[Patterns That Keep It Flexible]
\textbf{Strategy} for \texttt{SchedulingStrategy} (nearest car, least
load, zoning) so dispatch policy is a plug-in. \textbf{State} for each
elevator's motion (idle/up/down) to keep the movement logic clean.
Optionally \textbf{Observer} so displays update as a car moves. The car's
own sweep is the classic SCAN scheduling of its pending stops.
\end{ideabox}

\subsection*{Class design (C++ skeleton)}

\begin{lstlisting}
#include <bits/stdc++.h>
using namespace std;

enum class Direction { Up, Down, Idle };

class Elevator {
public:
    int id, floor = 0;
    Direction dir = Direction::Idle;
    set<int> targets;                           // pending stops

    void addTarget(int f) { targets.insert(f); }
    int distanceTo(int f) const { return abs(floor - f); }

    void step() {                               // one SCAN move
        if (targets.empty()) { dir = Direction::Idle; return; }
        int next = (dir == Direction::Down)
                     ? *prev(targets.upper_bound(floor))
                     : *targets.lower_bound(floor);
        if (next > floor) { floor++; dir = Direction::Up; }
        else if (next < floor) { floor--; dir = Direction::Down; }
        else { targets.erase(floor); cout << "Car " << id << " stop @ " << floor << "\n"; }
    }
};

struct SchedulingStrategy {                     // Strategy: which car serves
    virtual Elevator* choose(vector<Elevator>& cars, int floor) = 0;
    virtual ~SchedulingStrategy() = default;
};
struct NearestCar : SchedulingStrategy {
    Elevator* choose(vector<Elevator>& cars, int floor) override {
        Elevator* best = nullptr;
        for (auto& c : cars)
            if (!best || c.distanceTo(floor) < best->distanceTo(floor)) best = &c;
        return best;
    }
};

class Dispatcher {
    vector<Elevator> cars;
    unique_ptr<SchedulingStrategy> policy;
public:
    Dispatcher(int n, unique_ptr<SchedulingStrategy> p) : policy(move(p)) {
        for (int i = 0; i < n; i++) cars.push_back({i});
    }
    void hallCall(int floor) {
        Elevator* car = policy->choose(cars, floor);
        car->addTarget(floor);
        cout << "Assigned car " << car->id << " to floor " << floor << "\n";
    }
    void tick() { for (auto& c : cars) c.step(); }
};

int main() {
    Dispatcher d(2, make_unique<NearestCar>());
    d.hallCall(5);
    for (int i = 0; i < 6; i++) d.tick();
    return 0;
}
\end{lstlisting}

\begin{takeawaybox}
Elevator design shines when you separate \emph{dispatch} (which car, a
\textbf{Strategy}) from \emph{motion} (how a car sweeps its stops, a
\textbf{State}/SCAN algorithm). Discuss the scheduling goal explicitly.
Follow-ups: concurrency across cars, priority/express calls, and
starvation avoidance.
\end{takeawaybox}
